\documentclass{../Templates/sbs-exam}

\usepackage{ulem}
\usepackage{array}
\usepackage{enumitem}
\usepackage{tabularx}
\usepackage{ulem}
\usepackage{xcolor}
\usepackage{multirow}
\usepackage{multicol}
\usepackage{graphicx} % 图形支持包（必须）
\usepackage{fancyhdr} % 页眉控制包（必须）
\usepackage{titling} % 提供高级标题控制
\usepackage{lmodern} % 现代字体支持
\usepackage{amsmath, amssymb} % 数学符号和公式支持
\usepackage{lipsum} % 用于生成示例文本
\usepackage{caption} % 图表标题支持

% ===== 资源路径配置 =====
\graphicspath{
  {../Assets/}     % 公共资源
  {../Media/}      % 试卷专用资源
  {./midterm-g10-c1-fall-2025/} % 本试卷资源
}

% ===== 考试元数据配置 =====
\examtitle{Final Exam, 2025-2026 T1}
\examsubject{G10 C2 Math}{Jan 2026}{90 Minutes}{Shi Feng}
\MarksBreakdown{2 \text{ Marks}}{10}{2 \text{ Marks}}{20}{40 \text{ Marks}}{100 \text{ Marks}}

% 自定义命令用于处理记分内容
\newcommand{\imarks}[1]{\hfill \textbf{(#1 marks)}}
\newcommand{\totalmarks}[1]{\par\noindent\hfill \textbf{(Total for question = #1 marks)}\par}

\begin{document}

\maketitle

\examnotice

\makemarksheet

\examtools

\newpage

\makeatletter


\section*{Multiple Choice Questions (10 questions, 2 marks each)}
\begin{enumerate}
    \item What is the period of the function $y = \tan \theta$?
    \begin{enumerate}
        \item $90^\circ$
        \item $270^\circ$
        \item $360^\circ$
        \item $180^\circ$
    \end{enumerate}
    
    \item The derivative of $f(x) = x^3$ is:
    \begin{enumerate}
        \item $x^2$
        \item $3x^2$
        \item $3x$
        \item $x^3$
    \end{enumerate}
    
    \item The indefinite integral $\int x^2  dx$ equals:
    \begin{enumerate}
        \item $x^3 + C$
        \item $2x + C$
        \item $\frac{x^3}{3} + C$
        \item $\frac{x^2}{2} + C$
    \end{enumerate}
    
    \item Which of the following is the correct factorization of $x^2 - 5x + 6$?
    \begin{enumerate}
        \item $(x-2)(x-3)$
        \item $(x+2)(x+3)$
        \item $(x-1)(x-6)$
        \item $(x+1)(x+6)$
    \end{enumerate}
    
    \item The equation of a circle with center $(2, -3)$ and radius 5 is:
    \begin{enumerate}
        \item $(x+2)^2 + (y-3)^2 = 25$
        \item $(x-2)^2 + (y+3)^2 = 5$
        \item $(x+2)^2 + (y-3)^2 = 5$
        \item $(x-2)^2 + (y+3)^2 = 25$
    \end{enumerate}
    
    \item The value of $\log_{10} 100$ is:
    \begin{enumerate}
        \item 1
        \item 2
        \item 10
        \item 100
    \end{enumerate}
    
    \item The gradient of the curve $y = x^2$ at $x=1$ is:
    \begin{enumerate}
        \item 2
        \item 0
        \item 1
        \item 4
    \end{enumerate}
    
    \item The indefinite integral $\int 3x  dx$ equals:
    \begin{enumerate}
        \item $3x^2 + C$
        \item $3x + C$
        \item $\frac{3x^2}{2} + C$
        \item $\frac{x^2}{2} + C$
    \end{enumerate}
    
    \item The midpoint of the line segment joining points $(1, 2)$ and $(5, 8)$ is:
    \begin{enumerate}
        \item $(4, 5)$
        \item $(3, 5)$
        \item $(3, 4)$
        \item $(4, 4)$
    \end{enumerate}
    
    \item The function $y = e^x$ is:
    \begin{enumerate}
        \item Always increasing
        \item Always decreasing
        \item Constant
        \item Periodic
    \end{enumerate}
\end{enumerate}

\section*{Short Answer Questions (20 questions, 2 marks each)}

\subsection*{True/False}
State whether each statement is True or False.

\section*{True or False Questions (10 questions)}
\begin{enumerate}
    \item The graph of $y = \cos \theta$ has a period of $360^\circ$. \hfill True \quad False
    
    \item The derivative of a constant function is zero. \hfill True \quad False
    
    \item The integral of $x^n$ with respect to $x$ is $\frac{x^{n+1}}{n+1} + C$ for all real $n$. \hfill True \quad False
    
    \item The line $y = mx + c$ always has a gradient of $m$. \hfill True \quad False
    
    \item The circle equation $(x-a)^2 + (y-b)^2 = r^2$ has center $(a, b)$. \hfill True \quad False
    
    \item $\log_a (xy) = \log_a x + \log_a y$ for $a > 0$, $a \neq 1$. \hfill True \quad False
    
    \item The tangent to a curve at a point is perpendicular to the radius at that point. \hfill True \quad False
    
    \item The function $y = \ln x$ is defined for all real $x$. \hfill True \quad False
    
    \item For any positive number $a$, $a^0 = 1$. \hfill True \quad False
    
    \item The perpendicular bisector of a chord passes through the center of the circle. \hfill True \quad False
\end{enumerate}

\section*{Fill in the Blanks (10 questions)}
\begin{enumerate}
    \item The derivative of $f(x) = 5x^4$ is \underline{\hspace{2cm}}.
    
    \item The  integral $\int 3  dx$ equals \underline{\hspace{2cm}}.
    
    \item The solution to $\log_2 x = 3$ is $x =$ \underline{\hspace{2cm}}.
    
    \item The gradient of the line passing through $(0,0)$ and $(3,6)$ is \underline{\hspace{2cm}}.
    
    \item The value of $\sin 90^\circ$ is \underline{\hspace{2cm}}.
    
    \item The function $y = \tan \theta$ has vertical asymptotes at \underline{\hspace{2cm}} degrees.
    
    \item The circle with centre (2,1) and radius 4 has an equation of \underline{\hspace{2cm}}.
    
    \item The area of a triangle with base $b$ and height $h$ is \underline{\hspace{2cm}}.
    
   \item The distance between points $(2,3)$ and $(5,7)$ is \underline{\hspace{2cm}}.
    
    \item The exponential function $y = a^x$ has a base that must be \underline{\hspace{2cm}}.
\end{enumerate}

\newpage

\section*{Long Answer Questions}

\begin{enumerate}

\item % Q1
\begin{enumerate}
        \item[(a)] Find the equation of the tangent to the curve with equation
        \[
        y = \frac{1}{4}x^3 - 8x^{-\frac{1}{2}}
        \]
        at the point \( P(4, 12) \). Give your answer in the form \( ax + by + c = 0 \) where \( a \), \( b \) and \( c \) are integers. \imarks{5}
        
        The curve with equation \( y = f(x) \) also passes through the point \( P(4, 12) \). Given that
        \[
        f'(x) = \frac{1}{4}x^3 - 8x^{-\frac{1}{2}}
        \]
        
        \item[(b)] Find \( f(x) \) giving the coefficients in simplest form. \imarks{5}
    \end{enumerate}

\vspace{8cm}

\item % Q2

Given the function \( f(x) = (x - 3)(3x^2 + x + a) - 35 \) where \( a \) is a constant:
          \begin{enumerate}
              \item State the remainder when \( f(x) \) is divided by \( (x - 3) \). \imarks{2}
              \item Given that \( (3x - 2) \) is a factor of \( f(x) \), show that \( a = -17 \). \imarks{3}
              \item Using algebra and showing each step of your working, fully factorise \( f(x) \). \imarks{5}
          \end{enumerate}

\vspace{8cm}

\item % Q3
\begin{enumerate}
    \item The circle \(C_1\) has equation
    \[
    x^2 + y^2 + 8x - 10y = 29
    \]
    \begin{enumerate}
        \item[(i)] Find the coordinates of the centre of \(C_1\).
        \item[(ii)] Find the exact value of the radius of \(C_1\).
    \end{enumerate}\imarks{6}
    \item[(b)] The circle \(C_2\) has equation
    \[
    (x - 5)^2 + (y + 8)^2 = 52
    \]\\
    With detailed reasoning, prove that the circles \(C_1\) and \(C_2\) neither touch nor intersect.\imarks{4}
\end{enumerate}

\vspace{6cm}

\item % Q4


\begin{enumerate}
    \item[(a)] Given that
    \[
    2\log_3(x + 1) + \log_3(x - 1) = \log_3(2x + 3) + 1
    \]
    show that
    \[
    x^3 + x^2 - 7x - 10 = 0
    \] \imarks{5}

    \item[(b)] Given also that \(-2\) is a root of the equation
    \[
    x^3 + x^2 - 7x - 10 = 0
    \]
    \begin{enumerate}
        \item[(i)] Use algebra to find the other two roots of the equation. \vspace{0.5cm} \imarks{3}
        \item[(ii)] Hence solve
        \[
        2\log_3(x + 1) + \log_3(x - 1) = \log_3(2x + 3) + 1
        \] \imarks{2}
    \end{enumerate}
\end{enumerate}

\end{enumerate}

\end{document}